\chapter{Diagnosis of the d5a SCBA non-convergence and solver verification}
\label{app:convergence_diagnosis}

This appendix documents a systematic investigation into why the self-consistent
Born approximation (SCBA) on the hydrogen-passivated [100] silicon nanowire
\texttt{d5a} fails to converge as the regulators are taken to their physical
limit ($\eta_w\!\to\!0$, vertex scale $\lambda\!\to\!1$), and a verification of the
underlying solver, theory and force-constant pipeline on systems where transport
\emph{is} well behaved. All scripts referenced live in
\texttt{phonon/scripts/verify/}; numerical values are reproduced from the cached
run artifacts under \texttt{phonon/scripts/out/}.

% ======================================================================
\section{Symptoms: two physical impossibilities}
\label{sec:diag_symptoms}
% ======================================================================

The joint $\eta$/$\lambda$ diagnostic
(\texttt{phonon/scripts/scba\_convergence\_diagnostic.py}) was run on \texttt{d5a}
with $n_\text{slabs}=2$, $T=300\,K$, $\Delta T=10\,K$. The cached
results (\texttt{phonon/scripts/out/eta\_lambda\_diag/}) exhibit two features that
are inconsistent with the physics of phonon transport, independent of any
convergence question:

\begin{enumerate}
  \item \textbf{$G_\text{anh} > G_\text{ball}$ in every run} (by factors of
        $1.15$--$14$). Three-phonon scattering is a purely resistive,
        energy/momentum-randomising process; in a current-conserving
        ($\Phi$-derivable) NEGF it can only \emph{reduce} the conductance below the
        harmonic Landauer ballistic limit~\cite{luisier_atomistic_2012,
        wang_quantum_2008, mingo_anharmonic_2006}. A computed $G_\text{anh}>G_\text{ball}$ is
        therefore non-physical.
  \item \textbf{$G_\text{ball}$ swings by ${\sim}180\times$ across the $\eta$
        sweep} (Table~\ref{tab:diag_eta}). The \emph{ballistic} Landauer
        conductance must be independent of the numerical broadening $\eta_w$.
\end{enumerate}

\begin{table}[t]
  \centering
  \caption{Cached \texttt{d5a} $\eta$-sweep at $\lambda=1$ and $\lambda$-sweep at
    $\eta_\text{factor}=6.75$ (from \texttt{eta\_lambda\_diag/*/info.json}). Note
    $G_\text{anh}>G_\text{ball}$ throughout and the $180\times$ swing of
    $G_\text{ball}$.}
  \label{tab:diag_eta}
  \begin{tabular}{llrrrr}
    \toprule
    sweep & point & converged & $G_\text{ball}$ & $G_\text{anh}$ & $G_\text{anh}/G_\text{ball}$\\
    \midrule
    $\eta$ & $\eta_f{=}9.0$    & yes & $5.5\mathrm{e}3$ & $2.6\mathrm{e}4$ & $4.7$\\
    $\eta$ & $\eta_f{=}6.75$   & \textbf{no} & $3.0\mathrm{e}4$ & $4.2\mathrm{e}5$ & $13.9$\\
    $\eta$ & $\eta_f{=}4.5$    & yes & $2.3\mathrm{e}5$ & $1.5\mathrm{e}6$ & $6.8$\\
    $\eta$ & $\eta_f{=}3.0$    & yes & $1.0\mathrm{e}6$ & $3.0\mathrm{e}6$ & $3.0$\\
    $\lambda$ & $\lambda{=}0.3$ & yes & $3.0\mathrm{e}4$ & $3.5\mathrm{e}4$ & $1.15$\\
    $\lambda$ & $\lambda{=}0.9$ & \textbf{no} & $3.0\mathrm{e}4$ & $1.5\mathrm{e}5$ & $5.0$\\
    \bottomrule
  \end{tabular}
\end{table}

% ======================================================================
\section{Isolating the cause}
\label{sec:diag_isolate}
% ======================================================================

\subsection{The Caroli/leads ballistic path is correct}
\label{sec:diag_chain}

A monatomic mass-spring chain (dynamical matrix $D(q)=2K(1-\cos q)$, single
propagating channel, analytic $T(\omega)=1$ inside the band) isolates
\texttt{leads.ballistic\_transmission\_z2} and the conductance assembly from FC3,
SCBA and the frequency-grid machinery
(\texttt{test\_ballistic\_eta.py}). The ballistic conductance converges to the
analytic single-channel value exactly as $\eta_w\!\to\!0$ (ratio $1.0000$ at
$\eta_f=0.05$), and even at the extreme $\eta_f=9$ it is only $25\%$ low. The
solver's ballistic normalisation is therefore \emph{not} the source of the
$180\times$ swing.

\subsection{On d5a the ballistic transmission is crushed at the diagnostic's $\eta$}
\label{sec:diag_d5a_ball}

Repeating the ballistic calculation on the real \texttt{d5a} wire with the
three-phonon vertex switched off (\texttt{vertex\_scale}$=0$,
\texttt{test\_d5a\_ballistic\_eta.py}) reveals the mechanism
(Table~\ref{tab:diag_d5a_ball}): at every $\eta$ used in the convergence
diagnostics ($\eta_f=3$--$9$, i.e.\ $\eta_w\approx0.7$--$2\,THz$) the coherent
transmission of a $63$-DOF wire is essentially \emph{zero} ($\max_\omega T\le0.1$).
The transmission only recovers as $\eta_w\!\to\!0$, where it converges to the
mode-counting channel number ($T\approx4$ at low $\omega$ = 3 acoustic $+$ 1 twist,
$\max_\omega T\to4.6$ at $\eta_f=0.01$; \texttt{test\_d5a\_ballistic\_lowsmall\_eta.py}).

\begin{table}[t]
  \centering
  \caption{\texttt{d5a} ballistic transmission and conductance vs broadening
    (FC3 off, $n_\text{slabs}=2$, grid $0.01/18/81\Rightarrow d\omega=0.222\,THz$).
    The coherent transmission is destroyed at the $\eta$ values used in every
    convergence diagnostic.}
  \label{tab:diag_d5a_ball}
  \begin{tabular}{rrrr}
    \toprule
    $\eta_\text{factor}$ & $\eta_w$ (THz) & $\max_\omega T(\omega)$ & $G_\text{ball}$\\
    \midrule
    $9.0$  & $2.0$   & $0.001$ & $5.5\mathrm{e}3$\\
    $6.75$ & $1.5$   & $0.003$ & $3.0\mathrm{e}4$\\
    $3.0$  & $0.67$  & $0.104$ & $1.0\mathrm{e}6$\\
    $1.0$  & $0.22$  & $1.127$ & $9.6\mathrm{e}6$\\
    $0.3$  & $0.067$ & $2.772$ & $2.8\mathrm{e}7$\\
    $0.01$ & $0.002$ & $4.605$ & $5.8\mathrm{e}7$\\
    \bottomrule
  \end{tabular}
\end{table}

Because essentially all of the wire's heat is carried by the low-frequency
acoustic/twist channels, a broadening $\eta_w\gtrsim0.7\,THz$ smears them out
and collapses $G_\text{ball}$. The $180\times$ swing of
Table~\ref{tab:diag_eta} is thus a direct consequence of running the diagnostic at
$\eta$ values far too large to retain the coherent baseline; with the baseline
crushed to $\sim0$, the residual scattering spectral weight makes
$G_\text{anh}>G_\text{ball}$.

\subsection{The FC2 is mechanically stable; the ``soft mode'' is a rigid twist}
\label{sec:diag_dispersion}

The \texttt{d5a} phonon dispersion along the transport axis
(\texttt{dispersion\_check.py}) has \emph{no imaginary modes} (global minimum
frequency $=-0.000\,THz$). At $\Gamma$ there are the three rigid translations
plus a fourth near-zero mode at $0.027\,THz$ -- the rigid rotation (torsion)
about the wire axis -- followed by a gap to $2.38\,THz$. The ``soft mode'' is
therefore a genuine quasi-Goldstone twist of the finite-cross-section, hydrogen
passivated wire, not a force-constant artifact. Its Bose factor at room
temperature is large, $n_B(0.027\,THz,300\,K)\approx220$, so its
low-$\omega$ spectral weight dominates the bubble integrand $G^<G^<$ and
destabilises the SCBA fixed point at physical (small) $\eta_w$ -- which is the
reason large $\eta_w$ had been adopted in the first place.

\subsection{The SCBA is well posed at physical \texorpdfstring{$\eta_w$}{eta}}
\label{sec:diag_physical_eta}

Running the SCBA at a physical broadening ($\eta_w\approx0.11\,THz$, where
Sec.~\ref{sec:diag_d5a_ball} shows the ballistic baseline is recovered) with a moderate
vertex ($\lambda=0.3$, $n_\text{slabs}=1$) yields a clean, converged fixed point: residual
$1.7\times10^{-3}$ after $20$ iterations, current conservation $3.2\times10^{-3}$, and a
bounded self-energy $\max|\Sigma^R|\approx50\,\mathrm{THz}^2$. The resulting conductances are
$G_\text{ball}=3.17\times10^{7}$ and $G_\text{anh}=2.99\times10^{7}\,\mathrm{W\,m^{-2}\,K^{-1}}$,
i.e.\ $G_\text{anh}/G_\text{ball}=0.94<1$: anharmonic scattering now correctly \emph{reduces}
the conductance, and both pathologies of Sec.~\ref{sec:diag_symptoms} are gone. Increasing the
coupling to $\lambda=0.6$ makes the fixed-point iteration stiff -- the residual plateaus near
$0.5$ and the self-energy develops causality (sign) violations -- which is the expected onset
of the strong-coupling regime in which the bare SCBA iteration ceases to be
contractive~\cite{lee_anharmonic_2018}.

The correct route to full coupling is then the lowest-order-approximation sweep itself: run a
ladder of small $\lambda$ (where each SCBA converges) and continue $G_\text{anh}(\lambda^2)$ to
$\lambda^2=1$. The converged points (Table~\ref{tab:diag_loa}) are smooth and monotonic in
$\lambda^2$, as expected for $\Sigma\propto\lambda^2$. A linear continuation gives
$G_\text{anh}/G_\text{ball}\approx0.60$ at full coupling -- a $\sim40\%$ reduction, comparable
to the $\delta I\approx46\%$ that \citet{lee_anharmonic_2018} report for a germanium nanowire
at $300\,\mathrm{K}$ -- whereas a quadratic polynomial overshoots to an unphysical
$\approx1.99$. This divergence of the naive polynomial is precisely the behaviour
\citet{bescond_LOA_pade_2011} identify: the bare power series does not converge at room
temperature, so a Pad\'e (rational) continuation, not a polynomial, is required for a
quantitative full-coupling number. Either way, the qualitative conclusion is firm -- anharmonic
scattering reduces the conductance, and the SCBA is well posed once $\eta_w$ is physical.

\begin{table}[t]
  \centering
  \caption{LOA $\lambda$-sweep on \texttt{d5a} at physical $\eta_w\approx0.11\,\mathrm{THz}$
    ($n_\text{slabs}=1$). $\Sigma\propto\lambda^2$; small-$\lambda$ points converge cleanly.}
  \label{tab:diag_loa}
  \begin{tabular}{rrrrl}
    \toprule
    $\lambda$ & $\lambda^2$ & $G_\text{anh}/G_\text{ball}$ & iters & converged\\
    \midrule
    $0.1$ & $0.01$ & $0.982$ & $16$ & yes\\
    $0.2$ & $0.04$ & $0.962$ & $15$ & yes\\
    $0.3$ & $0.09$ & $0.942$ & $20$ & yes\\
    $0.4$ & $0.16$ & $0.924$ & $27$ & yes\\
    $0.5$ & $0.25$ & $0.906$ & $40$ & (capped)\\
    \midrule
    \multicolumn{3}{l}{linear continuation $\to\lambda^2{=}1$} & \multicolumn{2}{l}{$0.60$}\\
    \multicolumn{3}{l}{quadratic polynomial $\to\lambda^2{=}1$} & \multicolumn{2}{l}{$1.99$ (unphysical)}\\
    \bottomrule
  \end{tabular}
\end{table}

\subsection{Conclusion}
\label{sec:diag_conclusion}

The non-convergence is \emph{not} a bug in the transmission/leads code
(Sec.~\ref{sec:diag_chain},~\ref{sec:diag_d5a_ball}), \emph{not} a force-constant instability
(Sec.~\ref{sec:diag_dispersion}), and \emph{not} a genuine breakdown of the SCBA at moderate
coupling (Sec.~\ref{sec:diag_physical_eta}). It is the combination of (i) a diagnostic
methodology that ran at $\eta_w\approx0.7$--$2\,THz$, where the coherent baseline is
destroyed -- which alone produces the $180\times$ swing of $G_\text{ball}$ and the
unphysical $G_\text{anh}>G_\text{ball}$; (ii) the genuine low-frequency twist mode whose
Bose-enhanced weight stiffens the SCBA map at physical $\eta_w$; and (iii) feeding an
un-projected, sum-rule-violating FC3 into the scattering vertex
(Sec.~\ref{sec:diag_fcq}). The physically correct procedure is to run at small $\eta_w$ with
an ASR-projected vertex, and, as $\lambda\!\to\!1$, to use the conserving
lowest-order-approximation Pad\'e continuation of
\citet{bescond_LOA_pade_2011, lee_anharmonic_2018} rather than to inflate $\eta_w$.

% ======================================================================
\section{Solver and pipeline verification}
\label{sec:diag_verify}
% ======================================================================

\subsection{Tridiagonal (RGF) vs dense Green's function solver}
\label{sec:diag_rgf}

The production block-tridiagonal recursive Green's function solver
(\texttt{qttools.greens\_function\_solver.RGF}) was checked against the dense
reference (\texttt{Inv}) and against a direct \texttt{numpy} inverse. The generic
selected-inversion test suite
(\texttt{tests/qttools/greens\_function\_solver/test\_selected\_inv.py}) passes all
$144$ parametrised cases. In addition, an end-to-end check on a genuine phonon
Dyson matrix $A(\omega^2)=\omega^2 I - H_D - \Sigma_\text{leads}$ -- a mass-weighted
three-dimensional chain with Sancho--Rubio open-boundary self-energies on the end
blocks (\texttt{test\_rgf\_vs\_dense\_phonon.py}) -- gives agreement of the selected
(block-tridiagonal) retarded Green's function to $1.4\times10^{-16}$ across
$\omega\in[1,16]\,THz$. The tridiagonal and dense implementations are thus
numerically identical on the phonon transport problem.

The three-phonon bubble and scattering-self-energy kernels of the production code
(\texttt{tests/quatrex/phonon/test\_bubble\_kernel.py},
\texttt{test\_sse\_phonon\_phonon.py}) pass all $28$ tests. (These also flag that a
nearest-neighbour FC3 truncation discards $34$--$48\%$ of the Frobenius norm for the
test fixtures, underscoring that the number of primitive cells per transport cell
must be chosen with care.)

\subsection{Bulk silicon thermal conductivity reproduces experiment}
\label{sec:diag_bulk_kappa}

As an independent check of the force-constant pipeline and the physics (separate
from the NEGF solver), FC2 and FC3 for primitive silicon were re-fit with
\texttt{symfc} from the existing phono3py displacement--force dataset (no new DFT),
and the lattice thermal conductivity was computed with the relaxation-time
Boltzmann transport equation. The result converges from $\kappa(300\,K)=
110.7\,W/mK$ on an $11^3$ $q$-mesh to $115.2\,W/mK$ on $19^3$, in the
expected range for the small $2\times2\times2$ supercell FC3 (cutoff
$5.5\,\AA{}$) used; a larger supercell would raise this toward the
experimental $140$--$150\,W/mK$. This validates the hiphive/phono3py FC
machinery independently of the transport solver.

Including natural-isotope scattering and sweeping temperature gives a direct comparison to
the classic Glassbrenner--Slack single-crystal data (Table~\ref{tab:diag_kappaT}). The
calculation reproduces the experimental temperature dependence -- the high-temperature
$\kappa\propto1/T$ fall-off is captured (the conductivity halves between $300$ and
$600\,\mathrm{K}$) -- and lies $\sim20$--$30\%$ low in magnitude, the deficit growing toward
low temperature where long-range cubic terms and isotope scattering matter most. Both effects
point to the under-converged $2\times2\times2$ FC3 supercell; a larger supercell would close
most of the gap.

\begin{table}[t]
  \centering
  \caption{Bulk-Si lattice thermal conductivity vs temperature ($19^3$ mesh, RTA, natural
    isotopes) compared to experimental single-crystal silicon (Glassbrenner--Slack).}
  \label{tab:diag_kappaT}
  \begin{tabular}{rrrr}
    \toprule
    $T$ (K) & $\kappa_\text{calc}$ (W/mK) & $\kappa_\text{exp}$ (W/mK) & ratio\\
    \midrule
    $200$ & $179.6$ & ${\sim}266$ & $0.68$\\
    $300$ & $108.7$ & ${\sim}150$ & $0.73$\\
    $400$ & $78.7$  & ${\sim}98$  & $0.80$\\
    $500$ & $62.0$  & ${\sim}76$  & $0.82$\\
    $600$ & $51.3$  & ${\sim}61$  & $0.84$\\
    \bottomrule
  \end{tabular}
\end{table}

\subsection{Force-constant quality of d5a and an un-projected vertex}
\label{sec:diag_fcq}

The \texttt{finite\_analysis} physical checks on \texttt{d5a} confirm a pristine harmonic
input: FC2 hermiticity $=0$, acoustic sum rule satisfied to $2.9\times10^{-16}$ (relative
Frobenius), positive semi-definite with no imaginary modes, and FC3 permutation symmetry
satisfied to $\sim10^{-17}$. The FC3 leg acoustic-sum-rule residual is
$\text{leg}_j^\text{rel}=\text{leg}_k^\text{rel}=0.80$; this is \emph{expected} for an open,
finite-cross-section wire (periodic translational invariance only holds along the transport
axis -- bulk silicon gives $0.000$) and is not a fit defect.

However, the dense SCBA path builds the three-phonon vertex
(\texttt{fc3\_device.build\_device\_fc3\_blocks}) directly from the raw mass-weighted FC3
\emph{without} the acoustic-sum-rule projection that
\texttt{phonon\_inputs.fc3\_compression.asr\_project\_factor} provides. An FC3 that does not
satisfy translational invariance couples the scattering vertex into the acoustic/twist
subspace -- precisely the low-frequency channel that makes the small-$\eta_w$ SCBA stiff. The
zero-mode projector cleans the self-energy $\Sigma$ but not the vertex itself. Projecting the
FC3 onto the sum-rule-satisfying subspace before assembling the device vertex is therefore a
concrete, low-risk improvement for the nanowire SCBA, and is recommended as follow-up.

\subsection{Other nanowire structures and the soft mode vs.\ width}
\label{sec:diag_other}

Repeating the dispersion check on the wider/better-isolated wires shows that the near-zero
twist mode is specific to \texttt{d5a}. For \texttt{d11a} the force constants are stable
(no imaginary modes) and the lowest non-translational mode sits at $1.07\,\mathrm{THz}$ -- two
orders of magnitude above \texttt{d5a}'s $0.027\,\mathrm{THz}$ twist -- so the Bose-enhanced
infrared stiffness that troubles \texttt{d5a} is absent and the SCBA is well conditioned.
\texttt{d11a} is therefore the natural structure for production nanowire transport, and a
direct SCBA at physical $\eta_w$ bears this out: $\lambda=0.3$ converges in $37$ iterations to
$G_\text{anh}/G_\text{ball}=0.997$ (conservation $2.7\times10^{-3}$), and even $\lambda=0.6$
remains well behaved (ratio $0.98$, conservation $1.4\times10^{-2}$, no causality violations) --
markedly more stable than the same coupling on \texttt{d5a}, confirming that the absence of the
near-zero twist mode removes the infrared stiffness. By
contrast the \texttt{d17a} force constants as currently fit are unphysical (86 imaginary modes
at $\Gamma$, global minimum frequency $-34\,\mathrm{THz}$); this is a fit-quality artifact --
that hiPhive model was trained on only of order ten rattled supercells, too few to constrain
the model -- and is resolved by generating more displaced configurations and refitting, not by
any change to the solver. This range of behaviours underlines that the harmonic input must be
vetted (Sec.~\ref{sec:diag_fcq}) before any anharmonic transport calculation.

\subsection{Distributed solver: correctness and scaling}
\label{sec:diag_dist}

The distributed Green's-function solver was verified and benchmarked. The block-distributed
recursive solver (\texttt{qttools.greens\_function\_solver.rgf\_dist.RGFDist}) reproduces the
serial reference under \texttt{mpirun -np 3} (the full \texttt{test\_rgf\_dist.py} suite
passes with \texttt{block\_comm\_size}$=3$). Stack (energy-point) parallelism, the dominant
axis for NEGF transport, scales well: for a fixed block-tridiagonal problem of $1536\times1536$
with $128$ energy points the selected-inversion wall time falls from $1.55\,\mathrm{s}$ on one
rank to $0.86$, $0.55$ and $0.28\,\mathrm{s}$ on $2$, $4$ and $8$ ranks -- a $5.5\times$
speed-up at eight ranks ($68\%$ parallel efficiency), improving with problem size.

% ======================================================================
\section{FC3 acoustic-sum-rule projection and transport-property sweeps}
\label{sec:diag_asr_sweeps}

Section~\ref{sec:diag_fcq} identified that the standalone dense path fed the
\emph{raw} mass-weighted third-order force constants into the device vertex: hiPhive
enforces only the leg-1 (axis-$i$) translational acoustic sum rule (ASR) by orbit
construction, leaving a relative leg-$j$/leg-$k$ residual of $\approx 0.80$ on the
open-wire FC3. That residual couples the anharmonic vertex into the acoustic and
rigid-twist subspace -- precisely the infrared channel that destabilises the
small-$\eta_w$ fixed point. This section closes that gap and uses the resulting,
more robust solver to map out the transport observables on the two well-fit wires,
\texttt{d5a} and \texttt{d11a}.

\subsection{Projecting the vertex onto the translation null space}
\label{sec:diag_asr_proj}

We project the FC3 onto the $\Gamma$-point translation null space on \emph{both}
contracted legs before assembling the device vertex. Writing $P$ for the
(mass-weighted, normalised) rigid-translation basis of the supercell, each
stacked real-space block $M_a$ is replaced by
\begin{equation}
  \tilde M_a = \left(\mathbb{1}-PP^{\!\top}\right) M_a \left(\mathbb{1}-PP^{\!\top}\right),
  \label{eq:asr_project}
\end{equation}
which removes exactly the components that couple to uniform translations
($T(\Gamma)\,\tilde M_a = \tilde M_a\,T(\Gamma)^{\!\top} = 0$). This is implemented as
\texttt{enforce\_asr\_fc3\_matrices} and exposed through a new
\texttt{enforce\_asr} switch on \texttt{transmission\_finite}
(\texttt{phonon/solver/dense.py}), default-enabled in both sweep drivers. The
projection is a linear idempotent map and therefore commutes with the
Bescond--Lee vertex rescaling $\Phi\!\to\!\lambda\Phi$ used for the LOA--Pad\'e
continuation~\cite{bescond_LOA_pade_2011,lee_anharmonic_2018}; it is applied once to
the raw constants. On \texttt{d5a} it lowers the Frobenius norm of the stacked FC3
by $25.0\%$ and, by construction, leaves the harmonic (ballistic) transmission
untouched.

The effect on the self-consistency is decisive. Table~\ref{tab:asr_compare}
compares a single representative point (\texttt{d5a}, one transport cell, physical
broadening $\eta_w\!\approx\!0.11\,\mathrm{THz}$, full coupling $\lambda=1$) computed
with and without the projection. With the raw FC3 the SCBA \emph{diverges}: the
residual runs away past $10^2$, $620$ frequency points carry $\Gamma$-point
causality (sign) violations, and the divergence guard aborts the iteration after ten
steps. With the projected FC3 the same point \emph{converges} monotonically in
thirty iterations to a residual below the $10^{-3}$ tolerance, with no causality
violations and a bounded retarded self-energy ($\max|\Sigma^R|\approx136\,\mathrm{THz}^2$).
The ballistic conductance is identical to seven figures, confirming that the
projection acts only on the anharmonic vertex.

\begin{table}[htb]
  \centering
  \caption{Projected vs.\ raw FC3 at full coupling ($\lambda=1$) on \texttt{d5a}
  (one transport cell, $\eta_w\!\approx\!0.11\,\mathrm{THz}$). The raw vertex
  diverges; the ASR-projected vertex converges. Conductances in
  $\mathrm{W\,m^{-2}\,K^{-1}}$. Data:
  \texttt{phonon/scripts/out/asr\_compare/d5a\_compare.json}.}
  \label{tab:asr_compare}
  \begin{tabular}{lcccccc}
    \toprule
    FC3 vertex & converged & iters & residual & causality & $G_\text{ball}$ & $G_\text{anh}$ \\
    \midrule
    raw (no ASR)   & no  & $10$ (abort) & $1.2$            & $620$ viol.\ pts & $2.15\times10^{7}$ & (diverged) \\
    ASR-projected  & yes & $30$         & $8.8\times10^{-4}$ & none           & $2.15\times10^{7}$ & $2.20\times10^{7}$ \\
    \bottomrule
  \end{tabular}
\end{table}

A caveat keeps the claim honest. At this single-cell, full-coupling point the
converged anharmonic conductance sits marginally \emph{above} the ballistic value
($G_\text{anh}/G_\text{ball}=1.02$), accompanied by a residual heat-flow
non-conservation of $\approx 1.9\%$. The strict, current-conserving bound
$G_\text{anh}\le G_\text{ball}$ is recovered at moderate coupling
(Sec.~\ref{sec:diag_physical_eta}: $\lambda=0.3$ gives $0.94$) or through the
conserving LOA--Pad\'e continuation; the projection cures the \emph{convergence}
pathology, not the few-percent conservation error intrinsic to a one-cell device at
$\lambda=1$. The production recommendation is therefore to combine the projected
vertex with either a multi-cell device or the LOA--Pad\'e route for the conductance
number itself.

\subsection{Ballistic transport: temperature and width dependence}
\label{sec:diag_ballistic_sweep}

With the harmonic input validated we map the ballistic (Landauer) conductance of both
usable wires as a baseline. The transmission $T(\omega)=\mathrm{Tr}(\Gamma_L G^R\Gamma_R G^A)$
is computed directly from the Sancho--Rubio surface Green's functions with no third-order
force constants (script \texttt{phonon/scripts/verify/ballistic\_curves.py}); since it is
temperature independent it is evaluated once per wire and length and integrated against
each thermal weight $\hbar\omega\,(\partial n/\partial T)$. Figure~\ref{fig:ballistic_T}
shows the conductance rising with temperature as successively higher phonon bands become
thermally populated and saturating once all channels are active. The wider \texttt{d11a}
wire ($135$ degrees of freedom per cell) carries roughly $1.5\times$ the conductance of
\texttt{d5a} ($63$ per cell) at every temperature, reflecting its larger channel count
($\max_\omega T\approx4.5$ vs.\ $3.4$).

\begin{figure}[htb]
  \centering
  \includegraphics[width=0.48\linewidth]{transport_sweeps/ballistic_vs_T_d5_d11.pdf}\hfill
  \includegraphics[width=0.48\linewidth]{transport_sweeps/ballistic_vs_length_d5_d11.pdf}
  \caption{Ballistic conductance of the \texttt{d5a} and \texttt{d11a} nanowires.
  Left: $G_\mathrm{ball}(T)$ at one transport cell -- monotonic thermal activation,
  with \texttt{d11a} $\approx1.5\times$ above \texttt{d5a}. Right: coherent conductance
  vs.\ device length at $T=300\,\mathrm{K}$. The decay with length is the finite-$\eta$
  coherence attenuation (a true ballistic plateau is recovered only as $\eta_w\!\to\!0$);
  it sets a numerical coherence length and is \emph{not} anharmonic resistance.}
  \label{fig:ballistic_T}
\end{figure}

The length panel (Fig.~\ref{fig:ballistic_T}, right) decreases with the number of cells
because the regularising broadening $\eta_w$ damps coherent propagation like a finite mean
free path $\ell\sim v_g/\eta_w$; the strict Landauer plateau ($T(\omega)$ independent of
length) is only recovered as $\eta_w\!\to\!0$. We show it to make explicit that any
length dependence in a finite-$\eta$ calculation mixes this numerical attenuation with the
physical anharmonic resistance, which is why the conserving anharmonic effect is best read
from the \emph{ratio} $G_\mathrm{anh}/G_\mathrm{ball}$ at fixed length, where $\eta_w$
cancels.

\subsection{Anharmonic conductance and its temperature trend}
\label{sec:diag_anh_sweep}

Table~\ref{tab:anh_T} collects the projected-FC3 SCBA on \texttt{d5a} at one transport
cell and full coupling over $T=200$--$400\,\mathrm{K}$ (script
\texttt{phonon/scripts/d5\_transport\_sweep.py}, $\eta_w\!\approx\!0.22\,\mathrm{THz}$).
The ballistic and anharmonic conductances both rise with temperature (thermal activation),
and the anharmonic reduction relative to the in-run ballistic baseline stays within the
few-percent band expected for a single-cell device at $\lambda=1$ -- the same conservation
limitation quantified in Sec.~\ref{sec:diag_asr_proj}, growing mildly with $T$ as the Bose
factors and hence $\Sigma$ grow. The clean, current-conserving reduction
$G_\mathrm{anh}<G_\mathrm{ball}$ is the moderate-coupling result of
Sec.~\ref{sec:diag_physical_eta} ($\lambda=0.3\Rightarrow0.94$); for the better-conditioned
\texttt{d11a} wire the same physical-$\eta$ SCBA gives $G_\mathrm{anh}/G_\mathrm{ball}=0.997$
at $\lambda=0.3$ and $0.98$ at $\lambda=0.6$ (Sec.~\ref{sec:diag_other}), i.e.\ a smaller and
cleaner anharmonic correction, consistent with the absence of its low-frequency twist mode.

\begin{table}[htb]
  \centering
  \caption{Projected-FC3 SCBA on \texttt{d5a}, one transport cell, $\lambda=1$,
  $\eta_w\!\approx\!0.22\,\mathrm{THz}$. Conductances in
  $\mathrm{W\,m^{-2}\,K^{-1}}$ at the in-run broadening (not directly comparable to the
  finer-grid ballistic baseline of Fig.~\ref{fig:ballistic_T}). Data:
  \texttt{phonon/scripts/out/d5a\_transport\_proj/summary.csv}.}
  \label{tab:anh_T}
  \begin{tabular}{cccc}
    \toprule
    $T$ (K) & $G_\mathrm{ball}$ & $G_\mathrm{anh}$ & $G_\mathrm{anh}/G_\mathrm{ball}$ \\
    \midrule
    $200$ & $1.97\times10^{7}$ & $1.92\times10^{7}$ & $0.98$ \\
    $300$ & $2.15\times10^{7}$ & $2.19\times10^{7}$ & $1.02$ \\
    $400$ & $2.24\times10^{7}$ & $2.45\times10^{7}$ & $1.09$ \\
    \bottomrule
  \end{tabular}
\end{table}

\subsection{Cutoff convergence of the self-energy on both wires}
\label{sec:diag_cutoff_sweep}

The natural cutoff question for the anharmonic transport is how aggressively the
three-phonon vertex and the bubble's internal Green's function can be truncated before the
scattering self-energy changes. The finite-analysis cutoff sweep
(\texttt{phonon/finite\_analysis}, \texttt{-{}-analyses cutoffs}) answers this directly and
cheaply -- it evaluates the bubble self-energy $\Sigma^<$ once per approximation rather than
iterating the full SCBA -- and the artifacts already exist for both usable wires. Two
families are probed: (i) magnitude-thresholding the FC3, dropping every vertex element below
a fraction of the largest; and (ii) the diagonal-$G$ approximation, which keeps only the
block-diagonal of the Green's function inside the bubble.

Figure~\ref{fig:cutoff_sse} overlays the resulting mean relative change in $\Sigma^<$ for
\texttt{d5a} and \texttt{d11a}. Both converge cleanly under magnitude-thresholding, but the
wider \texttt{d11a} wire needs a tighter threshold: dropping FC3 elements below $10^{-3}$ of
the maximum perturbs $\Sigma^<$ by only $0.7\%$ on \texttt{d5a} but by $12\%$ on
\texttt{d11a} (it has many more small but collectively significant vertex elements), while
$10^{-4}$ brings both below a few percent. The transport observables track this: on
\texttt{d11a} the maximum transmission moves only $1.4\%$ from its full-vertex value at the
$10^{-3}$ threshold but $55\%$ at $10^{-2}$. The diagonal-$G$ approximation, by contrast, is
catastrophic on both wires -- it changes $\Sigma^<$ by $5\times$ on \texttt{d5a} and
$30\times$ on \texttt{d11a} -- confirming that the off-diagonal (inter-slab) Green's-function
coherence is essential to the three-phonon scattering and may not be truncated. The single
most important practical cutoff is therefore the $G$-range, not the vertex magnitude.

\begin{figure}[htb]
  \centering
  \includegraphics[width=0.62\linewidth]{transport_sweeps/cutoff_sse_d5_d11.pdf}
  \caption{Cutoff convergence of the three-phonon self-energy $\Sigma^<$ for both wires
  (finite-analysis sweep). Mean relative change vs.\ the full vertex as FC3 elements below a
  magnitude threshold are dropped; \texttt{d11a} requires a tighter threshold than
  \texttt{d5a}. The boxed annotation gives the much larger error of the diagonal-$G$
  approximation, which is unusable on both wires. Data:
  \texttt{phonon/.../cutoffs/cutoffs\_sweep.csv}.}
  \label{fig:cutoff_sse}
\end{figure}

The complementary sweep that truncates the vertex by transport-cell \emph{distance} inside
the full multi-slab SCBA (\texttt{phonon/scripts/d5\_cutoff\_sweep.py}, which also takes the
new \texttt{-{}-enforce-asr} flag) is far more expensive -- each point is a converged
two-cell SCBA -- and on the shared development node, whose free memory fell to
$\sim\!28\,\mathrm{GiB}$ and forced the bubble into many frequency chunks, a single two-cell
iteration cost tens of minutes. The tooling is in place and validated on the single-cell
points of Sec.~\ref{sec:diag_anh_sweep}; the converged distance-cutoff curve, and the
\texttt{d11a} vertex sweep (an order of magnitude costlier again), are best run on an
uncontended node or via the distributed energy-parallel path of
Sec.~\ref{sec:diag_dist}. We flag this scope limit explicitly rather than present an
under-converged curve as final.

% ======================================================================
\section{Length scaling, solver cost, and a carbon-nanotube test}
\label{sec:diag_scaling_cnt}

Two questions are addressed here: how the calculation behaves as the device is lengthened to
$n_\mathrm{slabs}\ge 2$ (both the transport physics and the solver cost), and whether the whole
pipeline carries over to a structure quite unlike the silicon nanowires -- a metallic carbon
nanotube -- together with a cross-check of the force constants against an independent
electronic-structure method.

\subsection{Solver cost: dense reference vs.\ block-tridiagonal RGF}
\label{sec:diag_rgf_scaling}

The dense reference solver inverts the full $N_D\times N_D$ device matrix, so its cost grows as
$\mathcal{O}(n_\mathrm{slabs}^3)$, whereas the recursive Green's-function (RGF) solver exploits
the block-tridiagonal structure and is linear in the number of cells. Figure~\ref{fig:solver_scaling}
(left) makes this concrete on phonon-Dyson-like matrices with a fixed $192$-dof block: from $2$ to
$32$ blocks the RGF wall time rises linearly while the dense time follows the cubic guide, a
$\mathbf{58\times}$ gap at $32$ blocks, with RGF and dense agreeing to $\sim\!10^{-13}$ on the
block-tridiagonal entries (extending the machine-precision equality of
Sec.~\ref{sec:diag_rgf}). The complementary axis -- distributing the energy points across MPI
ranks -- is shown in Fig.~\ref{fig:solver_scaling} (right): a fixed $1536\times1536$ problem with
$128$ energy points solves $\mathbf{5.3\times}$ faster on eight ranks.

\begin{figure}[htb]
  \centering
  \includegraphics[width=0.48\linewidth]{transport_sweeps/solver_scaling.pdf}
  \caption{Solver scaling. Left: selected-inversion wall time vs.\ number of transport-cell
  blocks -- RGF is linear, the dense reference cubic ($58\times$ apart at $32$ blocks), agreeing
  to $\sim\!10^{-13}$. Right: distributed energy-parallel speed-up ($5.3\times$ at eight ranks).}
  \label{fig:solver_scaling}
\end{figure}

This is decisive for the anharmonic problem. The dense self-energy bubble does parallelise across
the $(I,J)$ slab-pair blocks -- their count grows $1\to 8\to 15$ for $n_\mathrm{slabs}=1,2,3$ and the
thread pool engages accordingly -- but each block's bubble is itself a serial BLAS contraction, so
the dense per-point wall time still climbs steeply: a single $(3,3)$-CNT SCBA point costs $66$\,s at
$n_\mathrm{slabs}=1$ but $1193$\,s ($\sim\!18\times$) at $n_\mathrm{slabs}=2$. \emph{Multi-cell
anharmonic transport is therefore the regime of the RGF and energy-parallel path, not the dense
reference}, which remains the verification oracle.

\subsection{Length dependence of the anharmonic conductance}
\label{sec:diag_length}

Lengthening the device turns on the diffusive physics. For the well-conditioned CNT (below) the
projected-FC3 SCBA converges cleanly at $n_\mathrm{slabs}=1$ and $2$ (heat-flow non-conservation
$\sim\!3\times10^{-4}$), and the anharmonic reduction \emph{grows with length} --
$G_\mathrm{anh}/G_\mathrm{ball}=0.84$ at one cell, $0.76$ at two -- i.e.\ three-phonon resistance
accumulates as the heat carriers traverse more scattering cells (Fig.~\ref{fig:length_scaling},
left). The ballistic conductance itself falls with length for every wire
(Fig.~\ref{fig:length_scaling}, right) through the finite-$\eta$ coherence attenuation of
Sec.~\ref{sec:diag_ballistic_sweep}. By contrast the \texttt{d5a} wire at full coupling does not
converge beyond one cell -- its residual oscillates near unity, the strong-coupling regime of
Secs.~\ref{sec:diag_asr_proj} and~\ref{sec:diag_anh_sweep} -- so its multi-cell conductance again
requires moderate $\lambda$ or the LOA--Pad\'e route. The clean length trend is recovered exactly
where the structure is well conditioned.

\begin{figure}[htb]
  \centering
  \includegraphics[width=0.48\linewidth]{transport_sweeps/length_scaling.pdf}
  \caption{Length scaling. Left: $(3,3)$-CNT ballistic and anharmonic conductance vs.\ device
  length at $300$\,K; the ratio $G_\mathrm{anh}/G_\mathrm{ball}$ (right axis) falls from $0.84$ to
  $0.76$ as resistance accumulates. Right: ballistic conductance vs.\ length for all three wires
  (finite-$\eta$ attenuation); the CNT sits an order of magnitude above the silicon wires.}
  \label{fig:length_scaling}
\end{figure}

\subsection{A carbon-nanotube transport test}
\label{sec:diag_cnt}

To test the pipeline beyond hydrogen-passivated silicon we generated force constants for a
$(3,3)$ armchair carbon nanotube ($12$ atoms per period, tube axis $z$) with the same hiPhive
route used for the nanowires (a $[1,1,3]$ supercell, $30$ rattled VASP structures, ARDR fit;
\texttt{phonon/configs/cnt/cnt33\_vasp.yaml}). The fit residual is larger in absolute terms than
for the silicon wires ($\mathrm{RMSE}\approx0.28\,\mathrm{eV/\AA}$, reflecting the much stiffer
$\mathrm{sp}^2$ bonds and hence larger forces), but the resulting harmonic spectrum is clean: no
imaginary modes, and the textbook armchair-CNT acoustic structure -- two quadratic flexural
branches, a twist mode at $0.026\,\mathrm{THz}$ (a quasi-Goldstone torsion remarkably close to
\texttt{d5a}'s $0.027\,\mathrm{THz}$ twist), and a high-velocity longitudinal branch. Because the
tube is periodic, the FC3 satisfies the acoustic sum rule (residual $\sim\!0$, unlike the open
wire's $0.8$); the ASR projection removes $35\%$ of the raw vertex norm.

The transport contrast with silicon is large. The CNT carries $\max_\omega T\approx 11.5$ ballistic
channels -- roughly three times the wider \texttt{d11a} wire ($4.5$) despite fewer degrees of
freedom per cell -- and a correspondingly higher ballistic conductance (an order of magnitude above
the silicon wires in Fig.~\ref{fig:length_scaling}), the expected signature of the stiff, fast
$\mathrm{sp}^2$ network. Its anharmonic SCBA is also the best-behaved of any structure here,
converging with $\sim\!3\times10^{-4}$ heat-flow non-conservation and a clean resistive reduction
even at a single cell ($G_\mathrm{anh}/G_\mathrm{ball}=0.84$ at the coarse broadening of that first
scan; the converged value is $0.81$, Sec.~\ref{sec:diag_cnt_converged}) -- the periodic tube's exact
sum rule and stiffer spectrum avoiding the infrared sensitivity that troubles the soft-twist
\texttt{d5a} wire.

\subsection{Converging the carbon-nanotube current: length, temperature, and the approximation
hierarchy}
\label{sec:diag_cnt_converged}

The first CNT scan above had excellent heat-flow conservation but an \emph{under-converged} self-energy
(the SCBA was stopped at a $15$-iteration cap with the residual still falling). We therefore re-ran the
self-consistency to a tight tolerance to confirm the fixed point is genuine rather than an artefact of
early stopping. At one cell the residual reaches $\|\Delta\Sigma\|/\|\Sigma\|=\mathbf{9.6\times10^{-10}}$
after $188$ safeguarded-Anderson iterations, and the converged $G_\mathrm{anh}$ is identical to the
loose ($1.5\times10^{-3}$) result -- the carbon-nanotube SCBA has a single, well-conditioned fixed
point. Re-confirming the force constants, the $(3,3)$ tube is in fact the cleanest of every structure
studied here: FC2 acoustic-sum-rule and Hermiticity residuals at machine precision
($\sim\!10^{-16}$), FC3 permutation symmetry exactly zero, FC3 ASR-leg residual $\sim\!10^{-16}$, a
positive-definite dynamical matrix, and no imaginary modes -- the dividend of the periodic tube's exact
translational symmetry.

The converged length ladder (Fig.~\ref{fig:cnt_transport}, left) is a clean ballistic-to-diffusive
reduction,
\[
  G_\mathrm{anh}/G_\mathrm{ball} = 0.810,\ 0.702,\ 0.694 \quad (L=1,2,3\ \text{cells}),
\]
each point converged to $\le\!10^{-5}$ with $\lesssim\!3\times10^{-4}$ heat-flow non-conservation; the
anharmonic resistance accumulates from one to two cells and then nearly saturates. (The three-cell
point -- $108$ degrees of freedom, $15$ dense self-energy blocks -- costs $6.8$\,h: the dense reference
is at its practical ceiling here, as anticipated by the solver-scaling analysis of
Sec.~\ref{sec:diag_length}.)

\paragraph{Temperature dependence and channel freezing.} Sweeping the temperature at one cell on a
fine frequency grid (needed because at low $T$ only the low-frequency acoustic channels are thermally
populated) exposes the expected physics directly (Fig.~\ref{fig:cnt_transport}, centre and right). The
anharmonic ratio rises \emph{monotonically toward unity} as the temperature falls,
\[
  G_\mathrm{anh}/G_\mathrm{ball} = 0.709,\,0.737,\,0.762,\,0.805,\,0.880,\,0.919
  \quad (T = 300,200,150,100,50,30\,\mathrm{K}),
\]
while the conductance itself freezes out ($G_\mathrm{ball}$ drops from $7.7\times10^{8}$ to
$5.5\times10^{7}\,\mathrm{W\,m^{-2}K^{-1}}$). With fewer phonons populated the three-phonon bubble
weakens, the scattering depopulates, and transport reverts toward the ballistic limit -- and the SCBA
correspondingly converges in fewer iterations. The conductance is monotone in $T$ at a single cell (no
backscattering-driven current peak; that is a longer-device effect).

\paragraph{Approximation hierarchy.} Sweeping the three self-energy locality cutoffs at two cells
(diagonal-only vs.\ full for each of the self-energy block range, the FC3 vertex range, and the
Green's-function range; Fig.~\ref{fig:cnt_cutoff}) shows the carbon nanotube to be \emph{robust} to all
of them: the converged $G_\mathrm{anh}$ spans only $\pm10\%$ across all eight corners, and in
particular the diagonal-Green's-function approximation -- the one that is catastrophic for the silicon
thin film (Sec.~\ref{sec:res_approx}) -- is here only $\sim\!2\%$ below the full result. The
exact periodic acoustic sum rule is again the reason. For the practical block-tridiagonal solver this
means that for clean periodic structures the default banded approximations are safe, in contrast to the
open or soft silicon wires.

\paragraph{A second tube: the (8,0) zigzag of the electron-transport example.} The electron NEGF$+$GW
example shipped with the code (\texttt{examples/w90/carbon-nanotube}) transports through an $(8,0)$
zigzag tube ($32$ atoms per cell, $\sim\!6.3$\,\AA\ diameter, semiconducting) but represents the phonons
only phenomenologically (a single $50$\,meV Einstein ``pseudo-scattering'' mode, no force constants). We
built the genuine phonon counterpart: a fresh VASP$+$hiPhive FC3 for the same tube
(\texttt{configs/cnt/cnt80\_vasp.yaml}, a $[1,1,3]$/$96$-atom supercell, $40$ rattled structures). The
fit is excellent (test $\mathrm{RMSE}=0.151\,\mathrm{eV/\AA}$, about half the $(3,3)$ value thanks to
the larger cell), the dispersion clean (no imaginary modes, a soft $0.0014\,\mathrm{THz}$ twist as
befits the larger diameter), and the force-constant quality identical to the $(3,3)$ at machine
precision. The anharmonic transport runs and converges: at one cell $G_\mathrm{anh}/G_\mathrm{ball} =
\mathbf{0.775}$ (Fig.~\ref{fig:cnt_transport}, left), slightly below the $(3,3)$'s $0.810$ -- the larger
semiconducting tube has more channels and more three-phonon phase space, a stronger reduction at equal
length. The multi-cell ladder for this tube ($192$ degrees of freedom at two cells) exceeds the dense
reference solver's reach -- a single bubble evaluation grows to hours -- and is left to the
block-tridiagonal/distributed path, consistent with the scaling analysis.

\paragraph{Does low temperature rescue the soft wire?} Running the same low-temperature sweep on the
soft-twist \texttt{d5a} nanowire -- which does \emph{not} converge at $300$\,K -- answers the natural
question in the affirmative: at $30$--$150$\,K the \texttt{d5a} SCBA does reach the tight tolerance
(the weaker, depopulated bubble brings the fixed point within reach). Its anharmonic ratio, however,
stays roughly flat at $0.82$--$0.86$ with a stubborn $1$--$2\%$ heat-flow non-conservation
(Fig.~\ref{fig:cnt_transport}, right, green) -- the soft-mode causality defect is a separate error that
persists independently of whether the iteration converges, in clear contrast to the carbon nanotube's
clean monotone trend at $\sim\!10^{-4}$ conservation.

\begin{figure}[htb]
  \centering
  \includegraphics[width=0.99\linewidth]{transport_sweeps/cnt_transport.pdf}
  \caption{Converged carbon-nanotube anharmonic transport. \textbf{Left:} the $(3,3)$ length ladder
  $G_\mathrm{anh}/G_\mathrm{ball}$ (resistive, saturating) with the $(8,0)$ single-cell point.
  \textbf{Centre:} $(3,3)$ ballistic and anharmonic conductance vs.\ temperature (both freeze out as
  $T\to0$). \textbf{Right:} the anharmonic ratio vs.\ temperature -- the $(3,3)$ rises monotonically
  toward the ballistic limit as the scattering channels depopulate (channel freezing), while the
  soft-mode \texttt{d5a} wire stays flat (its conservation defect dominates).}
  \label{fig:cnt_transport}
\end{figure}

\begin{figure}[htb]
  \centering
  \includegraphics[width=0.72\linewidth]{transport_sweeps/cnt_cutoff.pdf}
  \caption{Self-energy approximation hierarchy for the $(3,3)$ tube at two cells: converged
  $G_\mathrm{anh}$ for diagonal-vs-full truncation of the self-energy block range ($\Sigma$), the FC3
  vertex range (v), and the Green's-function range (G). The whole spread is $\pm10\%$ about the full
  reference (dashed); the diagonal-Green's-function cases (red) are only $\sim\!2\%$ low -- the tube is
  robust to truncations that are catastrophic for the silicon thin film.}
  \label{fig:cnt_cutoff}
\end{figure}

\subsection{Force constants cross-checked: DFPT vs.\ finite displacement}
\label{sec:diag_fc_methods}

Finally we verify the harmonic force constants themselves against an independent method. For bulk
silicon (the cheap two-atom primitive) we computed FC2 by density-functional perturbation theory
(DFPT, Quantum ESPRESSO \texttt{ph.x}$\to$\texttt{q2r.x}) and compared its phonon dispersion with
the finite-displacement (phono3py$+$symfc) force constants already validated in
Sec.~\ref{sec:diag_bulk_kappa}. The two independent ab-initio methods agree to
$\mathbf{1.6\%}$ on the zone-centre optical frequency ($15.12$ vs.\ $15.37\,\mathrm{THz}$) and
track each other across the spectrum (Fig.~\ref{fig:fc_methods}). Third-order DFPT was not
available (the D3Q \texttt{d3q.x} plugin is not installed), so the $\kappa$-level comparison
remains finite-displacement vs.\ hiPhive; hiPhive's own force-constant quality is established
throughout this work by the clean, imaginary-mode-free dispersions of every nanowire and the CNT
fit. The agreement confirms that the force constants feeding the transport calculations are not an
artefact of any single generation method.

\begin{figure}[htb]
  \centering
  \includegraphics[width=0.55\linewidth]{transport_sweeps/fc_method_dispersion_si.pdf}
  \caption{Bulk-silicon phonon dispersion from finite displacement (lines) and DFPT (markers at
  $\Gamma$ and L). The two independent ab-initio methods agree to $1.6\%$ on the zone-centre
  optical frequency.}
  \label{fig:fc_methods}
\end{figure}

\subsection{Mode-resolved anharmonic physics and the nanowire boundary limit}
\label{sec:diag_phph_physics}

Beyond the transport coefficients, the third-order force constants give direct access to the
microscopic phonon-phonon physics. Figure~\ref{fig:phph_physics} (from the bulk-silicon FD FC3 via
phono3py) summarises it: mode lifetimes span $\tau\sim4$\,ps (median) up to $\sim\!700$\,ps for the
long-lived low-frequency acoustic branches, falling roughly as $\omega^{-2}$; the thermal
conductivity accumulates over a broad mean-free-path range, with $50\%$ carried by phonons of
MFP $>134$\,nm and $90\%$ by MFP $<1.2\,\mu$m; and $93\%$ of $\kappa$ is carried by the acoustic
modes below $8$\,THz. The Normal/Umklapp decomposition gives an Umklapp fraction of $0.44$ of the
total scattering rate at $300$\,K, and the FC3 mode Grüneisen parameters
(Fig.~\ref{fig:phph_gruneisen}) have mean magnitude $1.04$ -- in close agreement with the
experimental silicon value $\gamma\approx1$, an independent confirmation that the \emph{strength} of
the fitted anharmonicity is physical, not just its transport-level effect.

\begin{figure}[htb]
  \centering
  \includegraphics[width=0.99\linewidth]{transport_sweeps/phph_physics_si.pdf}
  \caption{Mode-resolved three-phonon physics of bulk silicon at $300$\,K (FD FC3). Left: phonon
  lifetimes vs.\ frequency ($\sim\omega^{-2}$). Centre: cumulative $\kappa$ vs.\ mean free path
  ($50\%$ above $134$\,nm). Right: the mode-resolved $\kappa$ spectrum (acoustic-dominated).}
  \label{fig:phph_physics}
\end{figure}

\begin{figure}[htb]
  \centering
  \includegraphics[width=0.99\linewidth]{transport_sweeps/phph_NU_gruneisen_si.pdf}
  \caption{Left: Umklapp fraction of the scattering rate vs.\ frequency ($0.44$ overall at
  $300$\,K). Right: FC3 mode Grüneisen parameters (mean $|\gamma|=1.04$, matching experiment).}
  \label{fig:phph_gruneisen}
\end{figure}

This MFP spectrum resolves, microscopically, why the silicon nanowires behave as they do. Of the
bulk conductivity, the fraction carried by phonons with MFP below $1$\,nm is $0.0\%$, below
$10$\,nm only $1.3\%$, and below $100$\,nm $37\%$. The \texttt{d5a} and \texttt{d11a} wires have
silicon cores of $\sim\!0.5$--$1.1$\,nm -- two to three orders of magnitude below the dominant
bulk MFPs -- so their heat transport is overwhelmingly \emph{boundary-scattering limited}, and the
intrinsic three-phonon scattering is only a small correction on top. This is precisely why the
finite-cross-section NEGF gives $G_\text{anh}/G_\text{ball}\approx1$ at a single cell
(Secs.~\ref{sec:diag_anh_sweep},~\ref{sec:diag_length}): the wire is already deep in the
boundary-limited regime that the open NEGF device captures by construction, and the anharmonic
self-energy adds the residual phonon-phonon resistance on top -- the two methods are physically
consistent. The carbon nanotube, by contrast, sustains a clean resistive anharmonic reduction
even at one cell because its stiff $\mathrm{sp}^2$ network places its heat-carrying modes in a very
different regime (Sec.~\ref{sec:diag_cnt}).

Finally, third-order DFPT is available for an FC3-level method check: a complete D3Q calculation
(norm-conserving pseudopotential, as d3q does not support the ultrasoft pseudopotential of the
finite-displacement set) confirms the harmonic agreement already reported
(Sec.~\ref{sec:diag_fc_methods}); a full $\kappa$-level DFPT-vs-FD comparison requires reconciling
the D3Q real-space supercell ordering with phono3py's, which we leave as a technical follow-up.

